% GENERATED FILE. Edit canonical lessons, then run npm run generate:books.
% canonical-source-hash: fnv1a64:512394454fdac017
% canonical-lessons: ML-C460-antarastra, ML-C460-agola, ML-C460-prayogikamaya, ML-C460-uddesikkuka, ML-C460-vyatyasappetuka, ML-R460-review-453-460, ML-R460-review-final-460

\chapter{Near and Far}
\label{ch:ml-near-and-far}
\hlchaptermodality{460}

\begin{chapteropening}
\textbf{By the end of this chapter:} \emph{I can say international, global, practical, to intend and to vary.}

Complete the last lesson of chapter 460: I can say international, global, practical, to intend and to vary.
\end{chapteropening}

\section[\texorpdfstring{\ml{അന്താരാഷ്ട്ര} (antārāṣṭra)}{അന്താരാഷ്ട്ര (antārāṣṭra)}]{\ml{അന്താരാഷ്ട്ര} (antārāṣṭra) — international}

\label{lesson:ML-C460-antarastra}



\begin{quote}
Before the new one: say the Malayalam for local, then the Malayalam for national.
\end{quote}

\subsection*{You'll want to know: \ml{അന്താരാഷ്ട്ര}}
\textbf{\ml{അന്താരാഷ്ട്ര}} — \emph{antārāṣṭra} — \textquotedblleft{}international\textquotedblright{}.

\begin{quote}
\ml{നഗരത്തിൽ} \ml{ഒരു} \ml{അന്താരാഷ്ട്ര} \ml{സമ്മേളനം} \ml{നടന്നു}.
\end{quote}

\begin{quote}
\emph{nagarattil oru antārāṣṭra sammēḷanaṁ naṭannu.}
\end{quote}

\begin{quote}
An international conference was held in the city.
\end{quote}

It stands before a noun: \ml{അന്താരാഷ്ട്ര} \ml{വിമാനത്താവളം}, an international airport.

\begin{cousinweb}
\textbf{\ml{അന്താരാഷ്ട്ര}} (\emph{antārāṣṭra}) is a modern word built from Sanskrit \emph{antar}, \textquotedblleft{}between, within\textquotedblright{} (CDIAL 356), and \emph{rāṣṭra}, \textquotedblleft{}a kingdom, a country\textquotedblright{} (CDIAL 10721): between countries.
\end{cousinweb}

\begin{grammarlens}[title={What you've built}]
One more word for saying what happens if.
\end{grammarlens}

\subsection*{Your turn}
\begin{itemize}
\raggedright
  \item \emph{Say it:} \emph{antārāṣṭra}
  \item \emph{Say it:} \emph{antārāṣṭra}, once more
  \item \emph{Say it:} \emph{nagarattil oru antārāṣṭra sammēḷanaṁ naṭannu.} — An international conference was held in the city.
  \item \emph{Recall:} say the Malayalam for local, then the Malayalam for national, then say \emph{antārāṣṭra} again
\end{itemize}

\begin{tcolorbox}[breakable,colback=teal!4,colframe=teal!35!black,title={Before you move on}]
What does \ml{അന്താരാഷ്ട്ര} mean? (\textquotedblleft{}International\textquotedblright{}.) Say it once more.
\end{tcolorbox}
\section[\texorpdfstring{\ml{ആഗോള} (āgōḷa)}{ആഗോള (āgōḷa)}]{\ml{ആഗോള} (āgōḷa) — global}

\label{lesson:ML-C460-agola}



\begin{quote}
Before the new one: say the Malayalam for national, then the Malayalam for international.
\end{quote}

\subsection*{You'll want to know: \ml{ആഗോള}}
\textbf{\ml{ആഗോള}} — \emph{āgōḷa} — \textquotedblleft{}global\textquotedblright{}.

\begin{quote}
\ml{ആഗോള} \ml{വിപണിയിൽ} \ml{എണ്ണയുടെ} \ml{വില} \ml{കുറഞ്ഞു}.
\end{quote}

\begin{quote}
\emph{āgōḷa vipaṇiyil eṇṇayuṭe vila kuṟaññu.}
\end{quote}

\begin{quote}
The price of oil fell on the global market.
\end{quote}

It stands before a noun: \ml{ആഗോള} \ml{പ്രശ്നം}, a global problem.

\begin{cousinweb}
\textbf{\ml{ആഗോള}} (\emph{āgōḷa}) is a modern learned word built from Sanskrit \emph{ā-}, \textquotedblleft{}up to, all round\textquotedblright{}, and \emph{gola}, \textquotedblleft{}a ball, a globe\textquotedblright{} (Monier-Williams): to do with the whole globe.
\end{cousinweb}

\begin{grammarlens}[title={What you've built}]
One more word for saying what happens if.
\end{grammarlens}

\subsection*{Your turn}
\begin{itemize}
\raggedright
  \item \emph{Say it:} \emph{āgōḷa}
  \item \emph{Say it:} \emph{āgōḷa}, once more
  \item \emph{Say it:} \emph{āgōḷa vipaṇiyil eṇṇayuṭe vila kuṟaññu.} — The price of oil fell on the global market.
  \item \emph{Recall:} say the Malayalam for national, then the Malayalam for international, then say \emph{āgōḷa} again
\end{itemize}

\begin{tcolorbox}[breakable,colback=teal!4,colframe=teal!35!black,title={Before you move on}]
What does \ml{ആഗോള} mean? (\textquotedblleft{}Global\textquotedblright{}.) Say it once more.
\end{tcolorbox}
\section[\texorpdfstring{\ml{പ്രായോഗികമായ} (prāyōgikamāya)}{പ്രായോഗികമായ (prāyōgikamāya)}]{\ml{പ്രായോഗികമായ} (prāyōgikamāya) — practical}

\label{lesson:ML-C460-prayogikamaya}



\begin{quote}
Before the new one: say the Malayalam for international, then the Malayalam for global.
\end{quote}

\subsection*{You'll want to know: \ml{പ്രായോഗികമായ}}
\textbf{\ml{പ്രായോഗികമായ}} — \emph{prāyōgikamāya} — \textquotedblleft{}practical\textquotedblright{}.

\begin{quote}
\ml{അവളുടെ} \ml{ആശയം} \ml{വളരെ} \ml{പ്രായോഗികമായിരുന്നു}.
\end{quote}

\begin{quote}
\emph{avaḷuṭe āśayaṁ vaḷare prāyōgikamāyirunnu.}
\end{quote}

\begin{quote}
Her idea was very practical.
\end{quote}

In practice is \ml{പ്രായോഗികമായി}.

\begin{cousinweb}
\textbf{\ml{പ്രായോഗികമായ}} (\emph{prāyōgikamāya}) is built from Sanskrit \emph{prāyogika}, \textquotedblleft{}applied, used, applicable\textquotedblright{} (Monier-Williams), from \emph{prayoga}, \textquotedblleft{}use\textquotedblright{}, with the adjective ending -\ml{ആയ}.
\end{cousinweb}

\begin{grammarlens}[title={What you've built}]
One more word for saying what happens if.
\end{grammarlens}

\subsection*{Your turn}
\begin{itemize}
\raggedright
  \item \emph{Say it:} \emph{prāyōgikamāya}
  \item \emph{Say it:} \emph{prāyōgikamāya}, once more
  \item \emph{Say it:} \emph{avaḷuṭe āśayaṁ vaḷare prāyōgikamāyirunnu.} — Her idea was very practical.
  \item \emph{Recall:} say the Malayalam for international, then the Malayalam for global, then say \emph{prāyōgikamāya} again
\end{itemize}

\begin{tcolorbox}[breakable,colback=teal!4,colframe=teal!35!black,title={Before you move on}]
What does \ml{പ്രായോഗികമായ} mean? (\textquotedblleft{}Practical\textquotedblright{}.) Say it once more.
\end{tcolorbox}
\section[\texorpdfstring{\ml{ഉദ്ദേശിക്കുക} (uddēśikkuka)}{ഉദ്ദേശിക്കുക (uddēśikkuka)}]{\ml{ഉദ്ദേശിക്കുക} (uddēśikkuka) — to intend; to mean}

\label{lesson:ML-C460-uddesikkuka}



\begin{quote}
Before the new one: say the Malayalam for global, then the Malayalam for practical.
\end{quote}

\subsection*{You'll want to know: \ml{ഉദ്ദേശിക്കുക}}
\textbf{\ml{ഉദ്ദേശിക്കുക}} — \emph{uddēśikkuka} — \textquotedblleft{}to intend; to mean\textquotedblright{}.

\begin{quote}
\ml{ഞാൻ} \ml{നിന്നെ} \ml{വേദനിപ്പിക്കാൻ} \ml{ഉദ്ദേശിച്ചില്ല}.
\end{quote}

\begin{quote}
\emph{ñān ninne vēdanippikkān uddēśiccilla.}
\end{quote}

\begin{quote}
I did not mean to hurt you.
\end{quote}

Past: \ml{ഉദ്ദേശിച്ചു}, \emph{uddēśiccu}.

What is meant takes -\ml{ആൻ}. \ml{ഉദ്ദേശ്യം}, a purpose, is the same family.

\begin{cousinweb}
\textbf{\ml{ഉദ്ദേശിക്കുക}} (\emph{uddēśikkuka}) is built from Sanskrit \emph{uddeśa}, \textquotedblleft{}pointing to, direction, aim\textquotedblright{} (Monier-Williams), with the Malayalam verb ending -\ml{ിക്കുക}.
\end{cousinweb}

\begin{grammarlens}[title={What you've built}]
One more word for saying what happens if.
\end{grammarlens}

\subsection*{Your turn}
\begin{itemize}
\raggedright
  \item \emph{Say it:} \emph{uddēśikkuka}
  \item \emph{Say it:} \emph{uddēśikkuka}, once more
  \item \emph{Say it:} \emph{ñān ninne vēdanippikkān uddēśiccilla.} — I did not mean to hurt you.
  \item \emph{Recall:} say the Malayalam for global, then the Malayalam for practical, then say \emph{uddēśikkuka} again
\end{itemize}

\begin{tcolorbox}[breakable,colback=teal!4,colframe=teal!35!black,title={Before you move on}]
What does \ml{ഉദ്ദേശിക്കുക} mean? (\textquotedblleft{}To intend; to mean\textquotedblright{}.) Say it once more.
\end{tcolorbox}
\section[\texorpdfstring{\ml{വ്യത്യാസപ്പെടുക} (vyatyāsappeṭuka)}{വ്യത്യാസപ്പെടുക (vyatyāsappeṭuka)}]{\ml{വ്യത്യാസപ്പെടുക} (vyatyāsappeṭuka) — to vary; to differ}

\label{lesson:ML-C460-vyatyasappetuka}



\begin{quote}
Before the new one: say the Malayalam for practical, then the Malayalam for to intend.
\end{quote}

\subsection*{You'll want to know: \ml{വ്യത്യാസപ്പെടുക}}
\textbf{\ml{വ്യത്യാസപ്പെടുക}} — \emph{vyatyāsappeṭuka} — \textquotedblleft{}to vary; to differ\textquotedblright{}.

\begin{quote}
\ml{ഓരോ} \ml{ജില്ലയിലും} \ml{വില} \ml{വ്യത്യാസപ്പെടുന്നു}.
\end{quote}

\begin{quote}
\emph{ōrō jillayiluṁ vila vyatyāsappeṭunnu.}
\end{quote}

\begin{quote}
Prices vary from district to district.
\end{quote}

Past: \ml{വ്യത്യാസപ്പെട്ടു}, \emph{vyatyāsappeṭṭu}.

The noun is \ml{വ്യത്യാസം}, a difference.

\begin{cousinweb}
\textbf{\ml{വ്യത്യാസപ്പെടുക}} (\emph{vyatyāsappeṭuka}) is built from Sanskrit \emph{vyatyāsa}, \textquotedblleft{}exchange; change, inversion\textquotedblright{} (Monier-Williams), with -\ml{പ്പെടുക} from native \ml{പെടുക} (DEDR 3853): to be changed about.
\end{cousinweb}

\begin{grammarlens}[title={What you've built}]
One more word for saying what happens if.
\end{grammarlens}

\subsection*{Your turn}
\begin{itemize}
\raggedright
  \item \emph{Say it:} \emph{vyatyāsappeṭuka}
  \item \emph{Say it:} \emph{vyatyāsappeṭuka}, once more
  \item \emph{Say it:} \emph{ōrō jillayiluṁ vila vyatyāsappeṭunnu.} — Prices vary from district to district.
  \item \emph{Recall:} say the Malayalam for practical, then the Malayalam for to intend, then say \emph{vyatyāsappeṭuka} again
\end{itemize}

\begin{tcolorbox}[breakable,colback=teal!4,colframe=teal!35!black,title={Before you move on}]
What does \ml{വ്യത്യാസപ്പെടുക} mean? (\textquotedblleft{}To vary; to differ\textquotedblright{}.) Say it once more.
\end{tcolorbox}
\section[(dialogue)]{Words for saying what happens if}

\label{lesson:ML-R460-review-453-460}



\begin{quote}
Say the Malayalam for to vote and to stand by.
\end{quote}

\subsection*{Your turn}
Say these aloud:

\begin{itemize}
\raggedright
  \item to vote, to stand by, an objection, to question, to respond
  \item to bring up, to defend, to uphold, to be convinced, to confirm
  \item to remind, to look up to, equal, ethical, to harm
  \item to include, to hold, to separate, to sort, to divide
  \item to appeal, to set in order, to carry out, to establish, to maintain
  \item to make use of, to allocate funds, to waste, a loan, to donate
  \item to impose, to collect dues, to manufacture, local, national
  \item international, global, practical, to intend, to vary
\end{itemize}

\begin{tcolorbox}[breakable,colback=teal!4,colframe=teal!35!black,title={Before you move on}]
To vote? (\textbf{\ml{വോട്ട്} \ml{ചെയ്യുക}}, \emph{vōṭṭŭ ceyyuka}.) To stand by? (\textbf{\ml{പിന്തുണയ്ക്കുക}}, \emph{pintuṇaykkuka}.) An objection? (\textbf{\ml{എതിർപ്പ്}}, \emph{etirppŭ}.) To question? (\textbf{\ml{ചോദ്യം} \ml{ചെയ്യുക}}, \emph{cōdyaṁ ceyyuka}.) To respond? (\textbf{\ml{പ്രതികരിക്കുക}}, \emph{pratikarikkuka}.) To bring up? (\textbf{\ml{ഉന്നയിക്കുക}}, \emph{unnayikkuka}.) To defend? (\textbf{\ml{പ്രതിരോധിക്കുക}}, \emph{pratirōdhikkuka}.) To uphold? (\textbf{\ml{ശരിവെക്കുക}}, \emph{śarivekkuka}.) To be convinced? (\textbf{\ml{ബോധ്യപ്പെടുക}}, \emph{bōdhyappeṭuka}.) To confirm? (\textbf{\ml{സ്ഥിരീകരിക്കുക}}, \emph{sthirīkarikkuka}.) To remind? (\textbf{\ml{ഓർമ്മിപ്പിക്കുക}}, \emph{ōrmmippikkuka}.) To look up to? (\textbf{\ml{ബഹുമാനിക്കുക}}, \emph{bahumānikkuka}.) Equal? (\textbf{\ml{തുല്യമായ}}, \emph{tulyamāya}.) Ethical? (\textbf{\ml{ധാർമ്മികമായ}}, \emph{dhārmmikamāya}.) To harm? (\textbf{\ml{ഉപദ്രവിക്കുക}}, \emph{upadravikkuka}.) To include? (\textbf{\ml{ഉൾപ്പെടുത്തുക}}, \emph{uḷppeṭuttuka}.) To hold? (\textbf{\ml{ഉൾക്കൊള്ളുക}}, \emph{uḷkkoḷḷuka}.) To separate? (\textbf{\ml{വേർതിരിക്കുക}}, \emph{vērtirikkuka}.) To sort? (\textbf{\ml{തരംതിരിക്കുക}}, \emph{taraṁtirikkuka}.) To divide? (\textbf{\ml{വിഭജിക്കുക}}, \emph{vibhajikkuka}.) To appeal? (\textbf{\ml{അഭ്യർത്ഥിക്കുക}}, \emph{abhyartthikkuka}.) To set in order? (\textbf{\ml{ക്രമീകരിക്കുക}}, \emph{kramīkarikkuka}.) To carry out? (\textbf{\ml{നിർവഹിക്കുക}}, \emph{nirvahikkuka}.) To establish? (\textbf{\ml{സ്ഥാപിക്കുക}}, \emph{sthāpikkuka}.) To maintain? (\textbf{\ml{നിലനിർത്തുക}}, \emph{nilanirttuka}.) To make use of? (\textbf{\ml{ഉപയോഗപ്പെടുത്തുക}}, \emph{upayōgappeṭuttuka}.) To allocate funds? (\textbf{\ml{വിനിയോഗിക്കുക}}, \emph{viniyōgikkuka}.) To waste? (\textbf{\ml{പാഴാക്കുക}}, \emph{pāḻākkuka}.) A loan? (\textbf{\ml{വായ്പ}}, \emph{vāypa}.) To donate? (\textbf{\ml{സംഭാവന} \ml{ചെയ്യുക}}, \emph{sambhāvana ceyyuka}.) To impose? (\textbf{\ml{ചുമത്തുക}}, \emph{cumattuka}.) To collect dues? (\textbf{\ml{പിരിക്കുക}}, \emph{pirikkuka}.) To manufacture? (\textbf{\ml{ഉത്പാദിപ്പിക്കുക}}, \emph{utpādippikkuka}.) Local? (\textbf{\ml{പ്രാദേശിക}}, \emph{prādēśika}.) National? (\textbf{\ml{ദേശീയ}}, \emph{dēśīya}.) International? (\textbf{\ml{അന്താരാഷ്ട്ര}}, \emph{antārāṣṭra}.) Global? (\textbf{\ml{ആഗോള}}, \emph{āgōḷa}.) Practical? (\textbf{\ml{പ്രായോഗികമായ}}, \emph{prāyōgikamāya}.) To intend? (\textbf{\ml{ഉദ്ദേശിക്കുക}}, \emph{uddēśikkuka}.) To vary? (\textbf{\ml{വ്യത്യാസപ്പെടുക}}, \emph{vyatyāsappeṭuka}.)
\end{tcolorbox}
\section[(dialogue)]{Near and Far — once more}

\label{lesson:ML-R460-review-final-460}



\begin{quote}
Say the Malayalam for international and global.
\end{quote}

\subsection*{Your turn}
\begin{itemize}
\raggedright
  \item \emph{Say it:} international, global, practical, to intend, to vary
\end{itemize}

\begin{tcolorbox}[breakable,colback=teal!4,colframe=teal!35!black,title={Before you move on}]
International? (\textbf{\ml{അന്താരാഷ്ട്ര}}, \emph{antārāṣṭra}.) Global? (\textbf{\ml{ആഗോള}}, \emph{āgōḷa}.) Practical? (\textbf{\ml{പ്രായോഗികമായ}}, \emph{prāyōgikamāya}.) To intend? (\textbf{\ml{ഉദ്ദേശിക്കുക}}, \emph{uddēśikkuka}.) To vary? (\textbf{\ml{വ്യത്യാസപ്പെടുക}}, \emph{vyatyāsappeṭuka}.)
\end{tcolorbox}
